% Generated draft: requires mathematical and source review.
\documentclass[11pt]{article}
\usepackage[margin=25mm]{geometry}
\usepackage{fontspec}
\setmainfont{Latin Modern Roman}
\usepackage{amsmath,amssymb,mathtools,unicode-math}
\setmathfont{Latin Modern Math}
\usepackage{enumitem,needspace,xurl,hyperref}
\setlength{\parindent}{0pt}
\setlength{\parskip}{5pt}
\setlength{\emergencystretch}{3em}
\newcommand{\N}{\mathbb{N}}
\newcommand{\Z}{\mathbb{Z}}
\newcommand{\Q}{\mathbb{Q}}
\newcommand{\R}{\mathbb{R}}
\newcommand{\C}{\mathbb{C}}
\newcommand{\F}{\mathbb{F}}
\newcommand{\A}{\mathbb{A}}
\newcommand{\PP}{\mathbb{P}}
\newcommand{\E}{\mathbb{E}}
\newcommand{\eps}{\varepsilon}
\newcommand{\dd}{\,\mathrm{d}}
\begin{document}
\section*{UnsolvedMath ID: 800735}
\textbf{Source problem code:} catalogue entry
\subsection{Definitions and mathematical statement}
A paratopological group is a group \(G\) equipped with a topology such that multiplication \(G\times G\to G\), \((x,y)\mapsto xy\), is jointly continuous. Continuity of inversion is not required. Write \(e\) for the identity and \(\mathcal N_G(e)\) for its neighborhoods. For subsets \(A,B\subseteq G\), put \(AB=\{ab:a\in A,\ b\in B\}\). Here countable includes finite.

The group \(G\) is \emph{left \(\omega\)-bounded} if
\[
\forall U\in\mathcal N_G(e)\ \exists F\subseteq G\quad
\bigl(F\text{ is countable and }G=FU\bigr).
\]
It is \emph{right \(\omega\)-bounded} if
\[
\forall U\in\mathcal N_G(e)\ \exists F\subseteq G\quad
\bigl(F\text{ is countable and }G=UF\bigr).
\]
It is \emph{\(\omega\)-bounded} if both properties hold. The countable set may depend on the neighborhood.

\textbf{Question (Guran; supplied source, Problem 9).} Does every right \(\omega\)-bounded paratopological group satisfy left \(\omega\)-boundedness? Explicitly, does the implication
\[
\begin{aligned}
&\bigl[\forall U\in\mathcal N_G(e)\ \exists F\subseteq G\quad
(F\text{ is countable and }G=UF)\bigr]\\
&\qquad\Longrightarrow
\bigl[\forall V\in\mathcal N_G(e)\ \exists H\subseteq G\quad
(H\text{ is countable and }G=HV)\bigr]
\end{aligned}
\]
hold for every paratopological group \(G\)?

The supplied excerpt specifies no separation axiom. Whether the cited source imposes a standing separation convention on paratopological groups remains unresolved; the intended class of groups requires confirmation.
\subsection{Short English statement}
Must a paratopological group in which every identity neighborhood covers the group by countably many right translates also admit a cover by countably many left translates of every identity neighborhood? The source's intended separation assumptions require confirmation.
\subsection{Sources}
\begin{itemize}
\item Local UnsolvedMath registry, permanent ID \(800735\), source code \texttt{catalogue entry}, titled \emph{Open Problems in Topological Algebra --- Problem 9 (Guran)}. The supplied record labels the problem open; this status is not independently verified.
\item Banakh and Choban, \emph{Some Open Problems in Topological Algebra} (2012), Problem 9, \url{https://arxiv.org/abs/1202.1619}. The supplied background attributes the problem to Guran; the registry lists Banakh and Choban in its proposer field.
\item Supplied provenance records access on 2026-10-16 and extraction by \texttt{arxiv\_latex}. These are supplied provenance claims, not evidence of independent source consultation for this draft. Supplied upstream revision: \texttt{d371ee917ac9647fbb06cd6fc529deec6f3dda82}.
\end{itemize}
\end{document}
